\documentclass[11pt, a4paper, twoside, titlepage]{scrbook}
\usepackage{ccicons}
\usepackage{fontawesome5} % Required for the cover watermark

% ==========================================
% KOMA-SCRIPT FONT SETTINGS & PREMIUM FONTS
% ==========================================
\usepackage{newpxtext, newpxmath}       % Palatino font for body and math
\usepackage[scaled=0.95]{inconsolata}   % Sleek modern monospace font for code
\usepackage[scaled=0.92]{helvet}       % Clean sans-serif for headings

\setkomafont{disposition}{\sffamily\bfseries}
\setkomafont{descriptionlabel}{\normalfont\bfseries}
\setkomafont{pageheadfoot}{\normalfont\small}

% ==========================================
% PACKAGES
% ==========================================
\usepackage[utf8]{inputenc}
\usepackage[T1]{fontenc}
\usepackage[english]{babel}
\usepackage{etoolbox}

% --- PREVENT \Bbbk AND \openbox COLLISIONS ---
\let\Bbbk\undefined
\let\openbox\undefined
% ---------------------------------------------

% Matematica e Informatica Teorica
\usepackage{amsmath, amssymb, amsthm, mathtools}
\usepackage{stmaryrd} % Simboli aggiuntivi

% Layout e Geometria
\usepackage{geometry}
\geometry{a4paper, inner=2.5cm, outer=2.5cm, bindingoffset=0.5cm, top=2.5cm, bottom=2.5cm}
\usepackage{multicol} % For 2-column acronyms

% Grafica e Colori
\usepackage{graphicx}
\usepackage{xcolor}
\usepackage{tikz}
\usetikzlibrary{automata, positioning, arrows}

% Box per Teoremi, Definizioni, Algoritmi
\usepackage{tcolorbox}
\tcbuselibrary{theorems, breakable, skins}

% Codice e Pseudocodice
\usepackage{listings}
\usepackage{algorithm}
\usepackage{algpseudocode}

% ==========================================
% CUSTOMIZATION & BRANDING (UNITN STYLE)
% ==========================================
\definecolor{unitnred}{RGB}{178, 35, 52}
\definecolor{codebg}{RGB}{245, 245, 245}

% ==========================================
% MODERN CHAPTER HEADINGS
% ==========================================
\makeatletter
\renewcommand*{\chapterlinesformat}[3]{%
  \Ifstr{#1}{chapter}{%
    \parbox[b]{\textwidth}{%
      \raggedright
      \Ifstr{#2}{}{}{% If there is a chapter number
        {\fontsize{70}{75}\selectfont\color{unitnred!20}#2}\\[-2ex]
      }%
      \Huge\bfseries\sffamily\color{black}#3\par\vspace{1cm}
    }%
  }{% For unnumbered chapters (like Preface/Appendix)
    \@hangfrom{#2}{#3}%
  }%
}
\makeatother

% ==========================================
% LINKS, REFERENCING & GLOSSARIES
% ==========================================
\usepackage{hyperref}
\hypersetup{
    colorlinks=true,
    linkcolor=unitnred,
    filecolor=magenta,      
    urlcolor=unitnred,
    citecolor=blue,
    pdftitle={Lecture Notes - Advanced Computing Architectures},
    pdfauthor={Luca Facchini},
    pdfsubject={Master's Degree in Computer Science - University of Trento},
}

\usepackage[acronym,toc]{glossaries-extra}
\usepackage{glossary-tree} % Defines base \subglossentry
\makeglossaries

% --- CUSTOM BULLETED SUB-ACRONYM STYLE ---
\newglossarystyle{acrobulletlist}{%
  \setglossarystyle{tree}%
  \renewcommand*{\subglossentry}[3]{%
    \par\vspace{2pt}%
    \hspace*{1.5em}\textbullet~\textbf{\glstarget{##2}{\glossentryname{##2}}} -- %
    \glossentrydesc{##2}\glspostdescription\space ##3\par
  }%
}
\setglossarystyle{acrobulletlist}

% --- AUTOMATED ACRONYM USAGE COUNTER HOOKS ---
\makeatletter
\appto\glslinkpostsetkeys{%
    \ifcsundef{acrocount@\glslabel}{%
        \csgdef{acrocount@\glslabel}{1}%
    }{%
        \csxdef{acrocount@\glslabel}{\the\numexpr\csuse{acrocount@\glslabel} + 1\relax}%
    }%
}

\renewcommand{\glsxtrpostdescacronym}{%
  \ifcsdef{acrocount@\glscurrententrylabel}{%
    \space\textit{[\csuse{acrocount@\glscurrententrylabel}]}%
  }{}%
}
\makeatother
\usepackage{array}

% CLEVEREF MUST ALWAYS BE LOADED LAST
\usepackage{cleveref}

% ==========================================
% ACRONYMS DEFINITIONS
% ==========================================
\newacronym{cmos}{CMOS}{Complementary Metal-Oxide-Semiconductor}
\newacronym{cpu}{CPU}{Central Processing Unit}
\newacronym{cpi}{CPI}{Clocks Per Instruction}
\newacronym{eembc}{EEMBC}{Embedded Microprocessor Benchmark Consortium}
\newacronym{gnd}{GND}{Ground}
\newacronym{ic}{IC}{Integrated Circuit}
\newacronym{ipc}{IPC}{Instructions Per Clock}
\newacronym{isa}{ISA}{Instruction Set Architecture}
\newacronym{mips}{MIPS}{Million Instructions Per Second}
\newacronym{mops}{MOPS}{Million Operations Per Second}
\newacronym{nmos}{nMOS}{N-type Metal-Oxide-Semiconductor}
\newacronym{nre}{NRE}{Non-Recurring Engineering}
\newacronym{pmos}{pMOS}{P-type Metal-Oxide-Semiconductor}
\newacronym{os}{OS}{Operating System}
\newacronym{spec}{SPEC}{Standard Performance Evaluation Corporation}
\newacronym{tpc}{TPC}{Transaction Processing Performance Council}
\newacronym{cisc}{CISC}{Complex Instruction Set Computer}
\newacronym{risc}{RISC}{Reduced Instruction Set Computer}
\newacronym{gpr}{GPR}{General-Purpose Register}
\newacronym{cu}{CU}{Control Unit}
\newacronym{simd}{SIMD}{Single Instruction, Multiple Data}
\newacronym{fp}{FP}{Floating-Point}
\newacronym[parent=fp]{fpu}{FPU}{\gls{fp} Unit}
\newacronym[parent=fp]{mflops}{MFLOPS}{Million \gls{fp} Operations Per Second}
\newacronym{alu}{ALU}{Arithmetic Logic Unit}
\newacronym{dsp}{DSP}{Digital Signal Processor}
\newacronym{cc}{CC}{Condition Code}
\newacronym{pc}{PC}{Program Counter}
\newacronym{agu}{AGU}{Address Generation Unit}
\newacronym{fifo}{FIFO}{First In First Out}
\newacronym{mac}{MAC}{Multiply-Accumulate}
\newacronym{fir}{FIR}{Finite Impulse Response}
\newacronym{iir}{IIR}{Infinite Impulse Response}
\newacronym{fft}{FFT}{Fast Fourier Transform}
\newacronym{bta}{BTA}{Branch Target Address}
\newacronym{ic-inst}{IC}{Instruction Count}
\newacronym{raw}{RAW}{Read After Write}
\newacronym{war}{WAR}{Write After Read}
\newacronym{waw}{WAW}{Write After Write}
\newacronym{ram}{RAM}{Random-Access Memory}
\newacronym[parent=ram]{dram}{DRAM}{Dynamic \gls{ram}}
\newacronym[parent=ram]{sram}{SRAM}{Static \gls{ram}}

% Stile del Codice
\lstset{
    backgroundcolor=\color{codebg},
    basicstyle=\ttfamily\footnotesize,
    keywordstyle=\color{unitnred}\bfseries,
    commentstyle=\color{green!60!black}\itshape,
    stringstyle=\color{blue},
    breaklines=true,
    frame=single,
    rulecolor=\color{lightgray},
    numbers=left,
    numberstyle=\tiny\color{gray},
    showstringspaces=false
}

% ==========================================
% TCOLORBOX THEOREMS & DEFINITIONS
% ==========================================
\newtcbtheorem[number within=chapter,
    crefname={definition}{definitions},
    Crefname={Definition}{Definitions},
    list inside=deflist
]{definition}{Definition}{%
    colback=unitnred!5,colframe=unitnred!80!black,fonttitle=\bfseries, breakable%
}{def}

\newtcbtheorem[number within=chapter,
    crefname={theorem}{theorems},
    Crefname={Theorem}{Theorems},
    list inside=thmlist
]{theorem}{Theorem}{%
    colback=blue!5,colframe=blue!80!black,fonttitle=\bfseries, breakable%
}{thm}

\newtcbtheorem[number within=chapter,
    crefname={lemma}{lemmas},
    Crefname={Lemma}{Lemmas},
    list inside=lemlist
]{lemma}{Lemma}{%
    colback=cyan!5,colframe=cyan!80!black,fonttitle=\bfseries, breakable%
}{lem}

\newtcbtheorem[number within=chapter,
    crefname={observation}{observations},
    Crefname={Observation}{Observations},
    list inside=obslist
]{observation}{Observation}{%
    colback=orange!5,colframe=orange!80!black,fonttitle=\bfseries, breakable%
}{obs}

\newtcbtheorem[number within=chapter,
    crefname={concept}{concepts},
    Crefname={Concept}{Concepts},
    list inside=cptlist
]{concept}{Concept}{%
    colback=green!5!white,colframe=green!60!black,fonttitle=\bfseries, breakable%
}{cpt}


% ==========================================
% HEADER & FOOTER (KOMA-Script Native)
% ==========================================
\usepackage[headsepline=0.4pt, footsepline=0.4pt, plainfootsepline]{scrlayer-scrpage}
\clearpairofpagestyles

\automark[section]{chapter}
\ohead{\headmark} 
\ihead{\textbf{Advanced Computing Architectures}} 

\ifoot[A.Y. 2026/2027 -- \textit{Notes by Luca Facchini}]{A.Y. 2026/2027 -- \textit{Notes by Luca Facchini}}
\ofoot[Page \pagemark]{Page \pagemark}

\addtokomafont{headsepline}{\color{unitnred}}
\addtokomafont{footsepline}{\color{unitnred}}
\setlength{\headheight}{24pt}

\newcommand{\lecturestop}[1]{%
    \par\vspace{1em}
    \begin{tcolorbox}[
        colback=gray!10,
        colframe=gray!60,
        coltext=gray!60!black,
        fontupper=\centering\bfseries\sffamily,
        boxrule=0.5pt,
        arc=2mm,
        left=4mm, right=4mm, top=2mm, bottom=2mm
    ]
        Lecture stopped here on #1 — Remaining slides to be covered in the next session.
    \end{tcolorbox}
    \par\vspace{1em}
}

% ==========================================
% DOCUMENT BEGINS
% ==========================================
\begin{document}
\frontmatter

% --- TITLE PAGE ---
\begin{titlepage}
    % Full-page background graphics using Ti*k*Z
    \begin{tikzpicture}[remember picture, overlay]
        % 1. Bold UniTn red vertical bar on the left edge
        \fill[unitnred] (current page.south west) rectangle ([xshift=1.5cm]current page.north west);
        % 2. Subtle gray geometric accent at the bottom right
        \fill[gray!10] (current page.south east) -- ([yshift=6cm]current page.south east) -- ([xshift=-8cm]current page.south east) -- cycle;
        % 3. Giant faint watermark icon in background
        \node[text=gray!10, font=\fontsize{250}{250}\selectfont, anchor=center] at (current page.center) {\faMicrochip};
    \end{tikzpicture}
    
    % Push content down and to the right to clear the red bar
    \vspace*{2cm}
    \noindent\hspace{1cm}%
    \begin{minipage}{\dimexpr\textwidth-1cm\relax}
        
        % University Logo
        \includegraphics[width=4.5cm]{UniTn.jpg} 
        
        \vspace{4cm}
        
        % Main Title Area
        {\Huge \bfseries \sffamily \color{unitnred} Advanced Computing Architectures\par}
        \vspace{0.5cm}
        {\LARGE \itshape Lecture Notes\par}
        \vspace{0.5cm}
        {\Large A.Y. 2026/2027\par}
        
        \vfill
        
        % Author Info Block
        \begin{flushleft}
            {\Large \textbf{Luca Facchini}}\\[0.2cm]
            \textit{Master's Degree in Computer Science}\\
            University of Trento\\[0.3cm]
            \href{mailto:luca.facchini-1@studenti.unitn.it}{\texttt{luca.facchini-1@studenti.unitn.it}}\\
            \href{mailto:luca@fc-software.it}{\texttt{luca@fc-software.it}}
        \end{flushleft}
        
        \vspace{1.5cm}
        
        % Footer line and date
        \textcolor{gray}{\rule{\linewidth}{0.8pt}}
        \vspace{0.2cm}\\
        {\small \textbf{Last updated:} \today \hfill \textit{University of Trento}}
        
    \end{minipage}
\end{titlepage}

% --- COPYRIGHT / LEGAL PAGE (Back of Title Page) ---
\newpage
\thispagestyle{empty}
\vspace*{\fill}

\noindent \textbf{About this Document}\\
These notes cover the topics of the \textbf{Advanced Computing Architectures} course, held by Prof. Passerone during the Master's Degree in Computer Science at the University of Trento. The content herein is derived from class lectures, course recordings, and the professor's official slides and notes.

\vspace{1em}
\noindent \textbf{License \& Copyright}\\
\textcopyright\ 2026 Luca Facchini. My personal contributions, transcriptions, and formatting within this document are distributed under the Creative Commons Attribution-NonCommercial 4.0 International license (\textbf{CC BY-NC 4.0} \ccbynceu). The full text of the license is available in the \texttt{LICENSE} file at the root of this project. Please note that any original intellectual property, figures, or direct excerpts sourced from the professor remain strictly under their respective original copyrights and licenses.

\vspace{1em}
\noindent \textbf{Disclaimer}\\
This material is provided ``AS IS'', without warranty of any kind, express or implied, including but not limited to the warranties of accuracy, completeness, or fitness for a particular purpose. These are independently produced student notes and do not constitute an official university publication. 

\vspace{1em}
\noindent \textit{If you find any typos or errors, please contact me at the email on the title page.}
\vspace{2cm}

\cleardoublepage

% --- TABLE OF CONTENTS ---
\addtokomafont{chapterentry}{\color{unitnred}}
\addtokomafont{chapterentrypagenumber}{\color{unitnred}}
\tableofcontents
\cleardoublepage
\mainmatter

% --- PREFACE / SCOPE ---
\addchap{Preface: Scope and Purpose}

\noindent \textbf{The Motivation}\\
These notes were initially born out of a personal need: the desire to organize, synthesize, and rigorously typeset the vast amount of information presented during the lectures. Writing and structuring this material in \LaTeX\ served as my primary method of active recall and study. However, recognizing how challenging it can be to study from scattered slides, recordings, and fragmented notes, I decided to open-source this project. The goal is to provide current and future students of the Master's Degree in Computer Science with a unified, structured, and accessible study companion.

\vspace{1em}
\noindent \textbf{The Scope}\\
The scope of this document is to cover the core theoretical concepts, definitions, mathematical proofs, and practical architectures discussed throughout the course. It is designed to act as a comprehensive bridge between the professor's official slides and the verbal explanations provided during class. 

\vspace{1em}
\noindent \textbf{How to Use This Material}\\
While I have made every effort to ensure the mathematical and technical rigor of this document, it is important to state what these notes are \textit{not}. They are \textbf{not a replacement} for the official recommended textbooks, nor are they a substitute for attending lectures and engaging with the professor. They are a student-made study guide. You should use them to clarify complex topics, review for exams, and quickly reference formal definitions.

\vspace{1em}
\noindent \textbf{Collaborative Spirit}\\
Because this is an actively maintained, open-source project, it is subject to continuous improvement. If you are studying from this document and spot a typographical error, a mathematical inaccuracy, or a missing concept, please do not hesitate to reach out via the contact information on the title page. Collaborative corrections are highly encouraged and deeply appreciated.

\vspace{2em}
\noindent \textit{Happy studying,}\\
\textbf{Luca Facchini}\\
BSc Graduate in Computer Science


\cleardoublepage

\chapter{Performance Metrics [14--16/09/2026]} 
This chapter explores the fundamental performance metrics, power scaling constraints, cost evaluation models, and quantitative design principles utilized in computer architecture.

\section{Power Consumption}
In modern \gls{cmos} technology, power consumption is driven by two primary components:
\begin{description}
    \item[Dynamic (Active)] Currents related to:
    \begin{itemize}
        \item The charging and discharging of logic gates during switching states.
        \item Transient short-circuit currents during signal transitions.
    \end{itemize}
    \item[Static] Leakage (sub-threshold) currents flowing continuously through transistors.
\end{description}
\noindent The prevailing historical trend shows that with each process generation, dynamic power scales by a factor of roughly $0.9$, while power density increases by a factor of $1.3$ to $1.8$, and static power consumption grows at an exponential rate.

\subsection{Dynamic Power Consumption}
\begin{definition}{Dynamic Power Consumption}{dyn-power}
    Dynamic power consumption is primarily caused by the charging and discharging of the capacitive load of every logic gate inside a digital \gls{cmos} integrated circuit:
    \begin{align*}
        \text{Energy/transition} &= C_L \cdot V_{DD}^2 \cdot P_{0\rightarrow1} \\
        \text{Power} &= \text{Energy/transition} \cdot f = C_{L} \cdot V_{DD}^2 \cdot P_{0\rightarrow1} \cdot f
    \end{align*} 
    Where:
    \begin{description}
        \item[\(C_L\) (Capacitance):] A function of wire length and transistor size.
        \item[\(V_{DD}\) (Supply Voltage):] Has steadily decreased across successive fabrication generations.
        \item[\(P_{0\rightarrow1}\) (Switching Probability):] The statistical probability of a logic transition.
        \item[\(f\) (Clock Frequency):] Higher frequencies increase the rate of gate state changes, raising power dissipation.
    \end{description}
\end{definition}

\begin{definition}{Short-Circuit Current}{scc}
    The finite slope of input signals causes a direct current path between \(V_{DD}\) and \gls{gnd} for a brief period during switching when both the \gls{nmos} and \gls{pmos} transistors are simultaneously conducting:
    \begin{align*}
        E_{sc} &= t_{sc} \cdot V_{DD} \cdot I_{\text{peak}} \cdot P_{0\rightarrow 1}\\
        P_{sc} &= t_{sc} \cdot V_{DD} \cdot I_{\text{peak}} \cdot f_{0\rightarrow 1}
    \end{align*}
    Where \(t_{sc}\) represents the rise/fall duration of the input signal, and \(I_{\text{peak}}\) is determined by the saturation current of the transistors and the input-to-output slope ratio.
\end{definition}

\subsection{Static Power Consumption}
\begin{definition}{Static Power Consumption}{stat-power}
    Sub-threshold current grows exponentially with an increase in temperature and a decrease in threshold voltage. As supply voltage scaling continues, sub-threshold conduction becomes the primary source of power dissipation:
    \[
        I_{\text{DSub}} = k \cdot e^{\frac{-q \cdot V_{T}}{a \cdot k_a \cdot T}}
    \]
    where \(k, a, k_a\) are technology-dependent constants, \(V_T\) is the threshold voltage, and \(T\) is the absolute temperature.
\end{definition}

\subsection{Total Power Consumption}
\begin{definition}{Total Power Consumption}{tot-power}
    Combining dynamic, short-circuit, and leakage components, the total power consumption of a system is expressed as:
    \[
        P_{\text{total}} = \left(C_{L} \cdot V^2_{DD} \cdot P_{0\rightarrow1} \cdot f\right) + \left(t_{sc} \cdot V_{DD} \cdot P_{0\rightarrow1} \cdot I_{\text{peak}} \cdot f\right) + \left(V_{DD} \cdot I_{\text{leakage}}\right)
    \]
    At the architectural level, designers manage these levers: \(f_{\text{clock}}\) (trading off performance), \(C_L\) (reducing gate counts), \(V_{DD}\) (bounded by transistor precision limits), and \(P_{0\rightarrow1}\) (optimized via microarchitectural design).
\end{definition}

\section{Cost Estimation}
Device cost tends to decrease over time even without improvements in fabrication technology, driven by the \textbf{learning curve effect}. By increasing production volumes, manufacturing costs drop significantly ($\sim 10\%$ cost reduction per doubling of volume), allowing non-recurring engineering (\gls{nre}) and development costs to be amortized over a massive machine base.

\begin{definition}{Integrated Circuit Cost}{ic-cost}
    The cost of a manufactured \gls{ic} is governed by:
    \[
        C = \frac{C_{\text{die}} + C_{\text{testing}} + C_{\text{packaging}}}{Y}
    \]
    where \(C_{\text{die}}\) is the manufacturing cost of the die, \(C_{\text{testing}}\) and \(C_{\text{packaging}}\) are fixed post-processing costs dependent on complexity, and \(Y\) is the final manufacturing yield.
\end{definition}

\begin{definition}{Die Cost and Wafer Capacity}{die-cost}
    Single-die cost depends directly on wafer yield and total available dies:
    \[
        C_{\text{die}} = \frac{C_{\text{wafer}}}{N_{\text{die}} \cdot Y_{\text{die}}}
    \]
    The total number of dies per wafer (\(N_{\text{die}}\)) is calculated accounting for wafer circular geometry and peripheral edge loss:
    \[
        N_{\text{die}} = \frac{\pi \cdot D^2}{4 \cdot A_{\text{die}}} - \frac{\pi \cdot D}{\sqrt{2 \cdot A_{\text{die}}}}
    \]
    This highlights the strong, non-linear cost penalty associated with larger die areas (\(A_{\text{die}}\)).
\end{definition}

\begin{definition}{Die Yield}{die-yield}
    Pre-packaging die yield is modeled statistically as:
    \[
        Y_{\text{die}} = Y_{\text{wafer}} \cdot \left(1 + \frac{N_{\text{defects}} \cdot A_{\text{die}}}{\alpha}\right)^{-\alpha}
    \]
    where \(N_{\text{defects}}\) is defect density per unit area, and \(\alpha\) ($\sim 4.0$) is an empirical parameter reflecting process complexity.
\end{definition}

\section{Performance Evaluation}
Historically, microprocessor performance scaling was driven by \textbf{Moore's Law} (\cref{cpt:moores-law}). Performance today scales faster than transistor density alone due to architectural parallelism.

\begin{concept}{Moore's Law}{moores-law}
    An empirical observation stating that the number of transistors in a dense \gls{ic} historically doubles approximately every 15 to 18 months.
\end{concept}

\begin{definition}{Ideal Properties of a Performance Metric}{perf_metric_properties}
    A robust performance metric must satisfy several key criteria: \textbf{Linearity} (proportional scale matching), \textbf{Reliability} (consistent relative ranking of systems \(X\) and \(Y\)), \textbf{Repeatability} (compatible variance across test runs), \textbf{Ease of Use}, \textbf{Consistency} (standardized definitions), and \textbf{Independence} (unambiguous mathematical formulation).
\end{definition}

\subsection{The Clock Rate}
While clock rate is repeatable, easy to measure, consistent, and independent, it is a **misleading metric** (pushed primarily for marketing). It is neither linear nor reliable because it ignores instructions per cycle, computational complexity per cycle, and memory/I/O bottlenecks.

\subsection{Instruction and Operation Metrics}
\begin{definition}{MIPS (Million Instructions Per Second)}{mips-def}
    Defined as the execution instruction throughput:
    \[
        \text{\gls{mips}} = \frac{N}{T \cdot 10^6}
    \]
    where \(N\) is total instruction count and \(T\) is execution time. While easy to measure, it is non-linear and inconsistent across different \glspl{isa}.
\end{definition}

\begin{definition}{MOPS and MFLOPS}{mops-mflops-def}
    Arithmetic throughput metrics measuring integer operations (\gls{mops}) and floating-point operations (\gls{mflops}) per second:
    \begin{align*}
        \text{\gls{mops}} &= \frac{O}{T \cdot 10^6} & \text{\gls{mflops}} &= \frac{F}{T \cdot 10^6}
    \end{align*}
    These metrics remain flawed because they fail to capture underlying hardware acceleration differences.
\end{definition}

\subsection{Execution Time}
\begin{definition}{Execution Time}{exec-time}
    The total wall-clock time required to complete a task on a computing device. Performance is strictly inversely proportional to execution time, making it the only metric that satisfies all ideal evaluation properties.
\end{definition}

\noindent \textbf{Execution Time vs. \gls{cpu} Time:} \gls{cpu} time measures pure processor engagement and is divided into \textbf{User time} (executing application instructions) and \textbf{System time} (executing \gls{os} kernel instructions on behalf of the task).

\begin{definition}{CPU Performance Equation}{cpu-perf-eq}
    Total \gls{cpu} time is evaluated using the fundamental performance equation:
    \[
        \text{\gls{cpu} Time} = N_{\text{CLK}} \cdot T_{\text{CLK}} = \text{CPI} \cdot \text{IC} \cdot T_{\text{CLK}}
    \]
    where \(\text{IC}\) is Instruction Count, \(\text{CPI}\) is \gls{cpi}, and \(T_{\text{CLK}}\) is the clock period.
\end{definition}

\begin{definition}{Average CPI}{avg-cpi}
    Computed across distinct instruction classes:
    \[
        \text{CPI} = \frac{\sum_{i=1}^n \text{IC}_i \cdot \text{CPI}_i}{\text{IC}}
    \]
\end{definition}

\subsection{Benchmarks and Normalization}
Architectural evaluations rely on standardized benchmark suites tailored to specific domains: Desktop Computing (\gls{spec}), Servers (\gls{tpc}), and Embedded Computing (\gls{eembc}).

\begin{definition}{Geometric Mean of Normalized Execution Times}{geom-mean}
    To aggregate benchmark performance without bias toward long or short tasks, suites like \gls{spec} compute the geometric mean of normalized execution times against a reference machine:
    \[
        G = \sqrt[N]{\prod_{j=1}^N \frac{T_{e_j}}{T_{R_j}}}
    \]
\end{definition}

\section{Quantitative Principles and System Optimization}
\begin{concept}{Architecture Optimization Goals}{arch-opt}
    System design is a multi-objective optimization problem balancing execution time (\(T\)), cost (\(C\)), and power (\(P\)). Desktop systems prioritize minimizing \(T \cdot C\) under a power budget (\(P \le P_{\text{max}}\)), whereas embedded systems prioritize minimizing \(C \cdot P\) under a hard deadline (\(T \le T_{\text{max}}\)).
\end{concept}

\begin{definition}{Pareto Point}{pareto-points}
    A design configuration is Pareto-optimal if no other configuration in the design space is better across all performance, cost, and power criteria.
\end{definition}

\subsection{Improving an Existing System}
\begin{definition}{Speedup}{speedup}
    The relative performance gain achieved by an architectural enhancement \(E\):
    \[
        \text{Speedup}(E) = \frac{\text{Execution Time without } E}{\text{Execution Time with } E}
    \]
\end{definition}

\begin{theorem}{Amdahl's Law}{amdahl-law}
    The performance improvement gained from using a faster execution mode is limited by the fraction of time that the faster mode can be used (\(F\)):
    \[
        S_{\text{overall}} = \frac{1}{(1 - F) + \frac{F}{S}}
    \]
    \begin{description}
        \item[Corollary 1:] Make the common case as fast as possible.
        \item[Corollary 2 (Law of Diminishing Returns):] Additional enhancements targeting the same fraction of computation yield diminishing overall returns.
    \end{description}
\end{theorem}

\lecturestop{16/09/2026}
\chapter{Instruction Set Architecture Principles}

\section{ISA Principles}
Fundamentally, an \gls{isa} is composed of an Instruction Set, an Instruction Classification, and Addressing Modes.

\begin{definition}{Instruction Set Architecture (\gls*{isa})}{isa}
    The \glsxtrfull{isa} includes all the rules describing the functional specification of a given hardware architecture. Each \gls{isa} is characterized by:
    \begin{itemize}
        \item Different classes of instruction operators;
        \item The number and type of operands;
        \item Supported addressing modes.
    \end{itemize}
\end{definition}

\noindent Architectures can be classified based on instruction complexity and count into \gls{cisc} (\textbf{Complex Instruction Set Computer}), which feature a lower ``semantic gap'' but require a highly complex \glsxtrfull{cu}, and \gls{risc} (\textbf{Reduced Instruction Set Computer}), which rely on a simpler, more modular organization at the cost of larger machine-level programs. Furthermore, \glspl{isa} can be distinguished by their internal storage topology: stack architectures, accumulator architectures, and \gls{gpr} (\textbf{General-Purpose Register}) architectures, as well as whether they implement a dedicated \texttt{load}/\texttt{store} paradigm or allow direct Register-Memory operations.

\paragraph{Instruction Definition} A machine-level instruction is a multiple-field binary string decoded by the processor's \gls{cu}, acting as the core bridge between hardware and software. Generally, an instruction is composed of an opcode, operand addresses, and a corresponding number of address qualifiers.

\noindent High-level instructions are typically classified into several major operator categories:
\begin{itemize}
    \item \textbf{Arithmetic:} Including basic arithmetic and logical operations;
    \item \textbf{Data Transfer:} Moving data between registers and memory;
    \item \textbf{Control:} Branching, jumping, and flow redirection;
    \item \textbf{System:} Privileged operations and hardware management;
    \item \textbf{\gls{fp}:} Specialized arithmetic for real numbers handled by an optional \gls{fpu} ISA extension;
    \item \textbf{Decimal and String:} Legacy primitives historically used for commercial data processing (largely deprecated in modern architectures);
    \item \textbf{Graphics and Multimedia:} \glsxtrfull{simd}-style instructions operating on multiple smaller data elements in parallel (e.g., executing four 16-bit operations within a single 64-bit operand).
\end{itemize}
While all computer architectures provide a full foundational set for Arithmetic, Data Transfer, and Control, system functions vary widely among architectures.

\subsubsection{Arithmetic Instructions}
Basic arithmetic and logical instructions are largely self-explanatory: \texttt{Add}, \texttt{Sub}, \texttt{Mul}, \texttt{And}, \texttt{Or}, and \texttt{Not}. 

\noindent Shift and rotate operations are typically categorized into:
\begin{itemize}
    \item \textbf{Arithmetic shifts:} Shift left (inserting zeros) and shift right (extending the sign bit to preserve numerical sign);
    \item \textbf{Logical shifts:} Shift left and shift right, filling all vacant bit positions uniformly with zeros;
    \item \textbf{Circular (rotate) shifts:} Bit rotation where bits shifted out on one end wrap around to the other (e.g., left rotate moves MSB into LSB; right rotate moves LSB into MSB).
\end{itemize}

\noindent Additionally, \textbf{comparison instructions} store no explicit destination result. Instead, they execute a subtraction or logical evaluation, and the \gls{alu} outputs update processor condition flags (such as Zero, Sign, Overflow, and Carry) to inform subsequent control flow decisions.

\subsubsection{Multimedia Instructions}
Multimedia operands are often narrower than the native 64-bit data word of modern processors. Rather than wasting 64-bit \gls{alu} capacity when operating on narrower bit-widths, multimedia instructions operate on several smaller data items in parallel. For example, a 64-bit register can be \textbf{partitioned} into four 16-bit segments, allowing a dedicated execution unit to perform four simultaneous additions in a single clock cycle. Similarly, this concept applies to floating-point operations. These parallel operations are formally classified as \glsxtrfull{simd} or vector instructions.

\subsubsection{Data Transfer Instructions} 
Data transfer instructions are dedicated commands that load data from main memory into registers, or store data from registers back to memory. This category also encompasses register-to-register moves and memory-to-memory transfers (though the latter is rarely implemented in \gls{risc} architectures). Finally, stack manipulation operations (pushing and popping) and hardware I/O communications through specific memory-mapped ports belong to this class.

\subsubsection{Control Flow Instructions}
To translate high-level constructs like \texttt{if} statements and \texttt{while} loops, the \gls{isa} must provide low-level equivalent instructions. These include unconditional jumps, conditional branches, subroutine calls, and returns. Furthermore, managing the temporary interruption of program execution is critical; this is achieved via \texttt{trap}, \texttt{exception}, and \texttt{wait} instructions. Additionally, some \glspl{dsp} feature specialized hardware \texttt{repeat} instructions for zero-overhead loop control.

\paragraph{Conditional Branches}
Most control flow is managed by conditional branches. For pipeline optimization, how the architecture evaluates the branch condition is a critical design choice. Generally, three approaches are used:

\begin{table}[!htbp]
    \centering
    \renewcommand{\arraystretch}{1.3} 
    \begin{tabular}{| >{\raggedright\arraybackslash}p{0.15\textwidth} | >{\raggedright\arraybackslash}p{0.12\textwidth} | >{\raggedright\arraybackslash}p{0.23\textwidth} | >{\raggedright\arraybackslash}p{0.15\textwidth} | >{\raggedright\arraybackslash}p{0.20\textwidth} |}
         \hline
         \textbf{Name} & \textbf{Examples} & \textbf{Enforcement} & \textbf{Advantages} & \textbf{Disadvantages}\\
         \hline\hline
         \textbf{\gls{cc}} & 80x86, ARM, PowerPC & Special bits are set by \gls{alu} operations, possibly under program control. & Sometimes conditions are already set implicitly by previous operations. & \gls{cc} introduces extra architectural state and constrains the ordering of instructions.\\
         \hline
         \textbf{Condition Register} & Alpha, MIPS & Tests evaluate an arbitrary register containing the boolean result of a comparison. & Simple hardware implementation and usage. & Consumes a general-purpose register.\\
         \hline
         \textbf{Compare and Branch} & PA-RISC, VAX & The comparison logic is integrated directly into the branch instruction itself. & Only one instruction is needed for a conditional branch. & May require too much work per instruction in heavily pipelined architectures. \\
         \hline
    \end{tabular}
    \caption{Comparison of Architectural Branch Condition Evaluation Strategies}
    \label{tab:branch-strategies}
\end{table}

\noindent Furthermore, since the vast majority of branch comparisons evaluate against zero, some architectures choose to treat comparisons-to-zero as a heavily optimized special case in hardware.

\subsubsection{System and Other Instructions}
All remaining non-standard instructions fall into this category. This includes, but is not limited to, virtual memory management, system calls, interrupt handling, peripheral management, complex bit manipulations, and iterative execution sequences (such as hardware division and square root operations).

\subsection{Why RISC?} 
Statistical profiling reported that across 5 representative programs from the \gls{spec} benchmark suite running on x86 architectures, \textbf{96\% of executed instructions fell within the 10 simplest instruction types}. 

In accordance with Amdahl's Law (\cref{thm:amdahl-law}), optimizing the common case yields vastly superior returns compared to over-engineering rare pathways. 

\begin{observation}{Amdahl's Law Application}{amdahl-application}
    Amdahl's law dictates that the overall performance improvement gained by optimizing a single part of a system is limited by the fraction of time that the targeted part is actually used. Consequently, focusing hardware optimization on the most frequent instructions provides maximum effectiveness.
\end{observation}

\noindent Consequently, \gls{risc} architectures are inherently more efficient. Modern hybrids (such as x86 processors) maintain internal hardware optimization layers that decode \gls{cisc} instructions dynamically into RISC-like micro-operations, balancing software complexity compatibility with high-performance execution cores.

\section{Addressing Modes}
Before detailing specific addressing modes, we must classify how architectures handle operand access. Specifically, instruction sets can be categorized by where they store operands and results.

\begin{concept}{Operand Storage Architectures}{operand-storage}
    Architectures define how operands are accessed and where results are stored. The four primary models are:
    \begin{description}
        \item[Stack:] Operands are implicitly pushed to and popped from the top of the stack.
        \item[Accumulator:] One operand is fetched from memory, while the other is implicitly located in the accumulator register, which also stores the final result.
        \item[Register-Memory:] One operand is fetched from memory, while the other is located in a register. The result is typically stored back into a register.
        \item[Register-Register (Load/Store):] Both operands must be loaded from memory into registers before execution. The \gls{alu} operates exclusively on registers, and the result is stored in a (potentially different) register.
    \end{description}
\end{concept}

\noindent Instructions typically feature 2 or 3 operands; in either case, one operand usually designates the destination location where the result shall be stored.

\begin{definition}{Addressing Modes}{addressing-modes}
    An addressing mode specifies how to calculate the effective memory address of an operand by using information held in registers and/or constants contained within a machine instruction. The primary modes include:
    \begin{description}
        \item[Immediate:] The operand is a constant explicitly specified within the instruction itself. Small data words are embedded directly into the instruction, while larger constants may be appended in the memory word immediately following the instruction.
        \item[Register:] The required operands are located in explicitly specified processor registers.
        \item[Register Indirect:] The operand is located in main memory at the address stored in a specified register. This allows for natural pointer manipulation.
        \begin{itemize}
            \item \textbf{Autoincrement Variant:} The pointer register's content is automatically updated (incremented or decremented) either before or after the memory access.
            \item \textbf{Displacement/Indexed Variant:} The effective address is calculated by adding a constant (displacement) or the value of another register (indexed) to the base pointer register.
        \end{itemize}
        \item[Memory Direct:] The operand is fetched from an explicit, absolute memory address provided directly within the instruction (or in the following word).
        \item[\gls{pc}-Relative:] A special case of displacement addressing where the base register is the \glsxtrfull{pc}. This is highly useful for calculating the destination addresses of branch instructions, which are typically known at compile time.
    \end{description}
\end{definition}

\section{Instruction Encoding}
Assembly language is a symbolic representation of the machine language of a specific processor, invented to abstract and simplify hardware-level programming. An assembly language source file fundamentally consists of two components: \textbf{instructions} and \textbf{directives}. 

The former are human-readable mnemonics of machine-level instructions that map directly to the processor's \gls{isa}. The latter (directives, sometimes called pseudo-operations) serve as meta-instructions for the assembler and the linker. Unlike assembly instructions, directives do not translate directly into executable machine code; rather, they control the assembly process, dictate memory alignment, allocate space for data, and organize the output binary structure.

\noindent Common assembly directives include:
\begin{description}
    \item[\texttt{.align n}] Forces the next instruction or data element to be loaded at the next available memory address where the lowest \(n\) bits are zero (effectively aligning memory to a \(2^n\) byte boundary).
    \item[\texttt{.ascii "string"}] Allocates sequential memory bytes to store the characters of the provided string.
    \item[\texttt{.byte}] Stores a list of explicitly provided 8-bit values into sequentially available memory.
    \item[\texttt{.data}] Instructs the assembler to place the subsequent items into the data segment (typically used for initialized global or static variables).
    \item[\texttt{.double} | \texttt{.int}] Allocates memory and sequentially stores the provided numbers as double-precision \gls{fp} or integer values, respectively.
    \item[\texttt{.global label}] Exports the specified label, making it globally visible to the linker and available for cross-referencing by code in other object files.
    \item[\texttt{.text}] Instructs the assembler to place the subsequent instructions into the text segment (the memory area strictly reserved for executable code).
\end{description}

Instruction encoding fundamentally affects two aspects of a computing system: the size of the compiled programs and the complexity of the processor's hardware implementation. To address varying architectural needs, three primary encoding schemes are utilized:
\begin{description}
    \item[Variable:] Suitable for architectures that support many addressing modes. The \texttt{opcode} length and the number of operands can vary, minimizing program size but significantly complicating the \gls{cu} decoder.
    \item[Fixed:] The \texttt{opcode} and the number of operands have a fixed length. This is ideal for architectures with few addressing modes and simplifies hardware decoding, though it typically leads to larger overall program sizes.
    \item[Hybrid:] The format length is dictated by the \texttt{opcode}. It combines a fixed-length base with one or two optional extension fields, balancing code density with decoding simplicity.
\end{description}

\paragraph{DLX Instruction Set}
DLX instructions are strictly fixed at 32-bits wide with a 6-bit opcode. The architecture allows for only three formats (Immediate, Register, and Jump) and explicitly supports only three addressing modes: Register, Immediate, and Displacement.

\begin{figure}[!htbp]
    \centering
    \begin{tikzpicture}[
        % Base style for instruction fields
        cell/.style={draw, rectangle, minimum height=8mm, align=center, fill=codebg, font=\sffamily\footnotesize},
        bitlabel/.style={font=\sffamily\tiny\color{gray}},
        typelabel/.style={font=\sffamily\bfseries, anchor=east}
    ]
        
        % ---------------------------------------------------
        % R-TYPE INSTRUCTION
        % ---------------------------------------------------
        \node[typelabel] at (-0.5, 0) {R-Type};
        
        % 6 bits = 2.4cm, 5 bits = 2.0cm
        \node[cell, minimum width=2.4cm, anchor=west] (R_op) at (0,0) {Opcode};
        \node[cell, minimum width=2.0cm, anchor=west] (R_rs) at (2.4,0) {rs};
        \node[cell, minimum width=2.0cm, anchor=west] (R_rt) at (4.4,0) {rt};
        \node[cell, minimum width=2.0cm, anchor=west] (R_rd) at (6.4,0) {rd};
        \node[cell, minimum width=2.0cm, anchor=west] (R_sh) at (8.4,0) {shamt};
        \node[cell, minimum width=2.4cm, anchor=west] (R_fn) at (10.4,0) {funct};
        
        % Bit length labels
        \node[bitlabel, above=1mm of R_op] {6 bits};
        \node[bitlabel, above=1mm of R_rs] {5 bits};
        \node[bitlabel, above=1mm of R_rt] {5 bits};
        \node[bitlabel, above=1mm of R_rd] {5 bits};
        \node[bitlabel, above=1mm of R_sh] {5 bits};
        \node[bitlabel, above=1mm of R_fn] {6 bits};

        % ---------------------------------------------------
        % I-TYPE INSTRUCTION
        % ---------------------------------------------------
        \node[typelabel] at (-0.5, -1.5) {I-Type};
        
        % 16 bits = 6.4cm
        \node[cell, minimum width=2.4cm, anchor=west] (I_op) at (0,-1.5) {Opcode};
        \node[cell, minimum width=2.0cm, anchor=west] (I_rs) at (2.4,-1.5) {rs};
        \node[cell, minimum width=2.0cm, anchor=west] (I_rt) at (4.4,-1.5) {rt};
        \node[cell, minimum width=6.4cm, anchor=west] (I_im) at (6.4,-1.5) {Immediate};
        
        % Bit length labels
        \node[bitlabel, above=1mm of I_im] {16 bits};

        % ---------------------------------------------------
        % J-TYPE INSTRUCTION
        % ---------------------------------------------------
        \node[typelabel] at (-0.5, -3.0) {J-Type};
        
        % 26 bits = 10.4cm
        \node[cell, minimum width=2.4cm, anchor=west] (J_op) at (0,-3.0) {Opcode};
        \node[cell, minimum width=10.4cm, anchor=west] (J_of) at (2.4,-3.0) {PC Offset};
        
        % Bit length labels
        \node[bitlabel, above=1mm of J_of] {26 bits};

    \end{tikzpicture}
    \caption{DLX 32-bit Instruction Encoding Formats}
    \label{fig:dlx-encoding}
\end{figure}

\subsection{Specific \glsfmtshort{isa} features in \glsfmtfullpl{dsp}}
Digital signal processors are widely used in several embedded applications, including wireless communications, consumer audio, hard disk drives, multimedia, instrumentation, and the automotive industry. Compared to other embedded applications, \glspl{dsp} are characterized by: computationally demanding numeric algorithms (often relying on vector dot products), high sensitivity to small numeric errors, strict real-time requirements, continuous analog-to-digital data streams, high data bandwidth, and highly predictable memory access and program flow. For these reasons, a dedicated \gls{isa} for \glspl{dsp} must be highly optimized for mathematical calculations and support complex addressing modes.

\paragraph{\glsxtrfull{mac}} 
The sum of products is the foundational operation for most digital signal processing algorithms:
\[
\begin{aligned}
    \text{FIR Filter} &\qquad y(n) = \sum_{k=0}^M a_k x(n-k)\\
    \text{IIR Filter} &\qquad y(n) = \sum_{k=0}^M a_k x(n-k) + \sum_{k=1}^N b_k y(n-k)\\
    \text{Convolution} &\qquad y(n) = \sum_{k=0}^N x(k) h(n-k)\\
    \dots&\qquad\dots
\end{aligned}
\]
In accordance with Amdahl's law (\cref{thm:amdahl-law}), since these calculations constitute the bulk of execution time, all modern \glspl{dsp} are provisioned with a dedicated hardware Multiply-Accumulate unit and a corresponding \gls{mac} instruction in their \glspl{isa}.

\paragraph{Zero-Overhead Looping (\texttt{Repeat})} 
\gls{dsp} applications require the repeated execution of a small number of arithmetic or multiplication instructions. In general-purpose processors, software looping uses branch and compare instructions to jump back to the beginning of the loop, and a loop counter must be actively maintained and incremented. Both of these actions require an extra register and consume data path clock cycles. To circumvent this, \glspl{dsp} provide specific hardware \texttt{repeat} instructions that execute loops with zero instruction-fetch overhead.

\subsubsection{Addressing and Memory Requirements} 
Most \gls{dsp} algorithms require access to large amounts of (usually sequential) data. They need the ability to perform multiple memory accesses within a single instruction cycle, and the ability to rapidly calculate new addresses in parallel with ongoing arithmetic operations. This is resolved through a modified Harvard (or Super-Harvard) dual-memory architecture, a large variety of flexible register-indirect addressing modes, and a dedicated functional unit devoted exclusively to calculating data addresses: the \glsxtrfull{agu}.

\paragraph{The \glsfmtshort{agu}}
The \gls{agu} is a dedicated piece of hardware consisting of a specific \gls{alu} and separate registers intended strictly for calculating memory addresses. Thus, address generation for subsequent instructions can be executed strictly in parallel with the main \gls{alu} datapath.

% AGU Schema Placeholder

\paragraph{Special Addressing Mode: Circular}
As previously discussed, \glspl{dsp} deal with iterative computations on continuous sets of data. \glsxtrfull{fifo} buffers and circular buffers are used to store data in adjacent memory locations. Normally, a program must maintain one read pointer and one write pointer, checking at the end of each operation if the buffer boundary has been reached.

\noindent To optimize this, the \gls{agu} contains special registers (often denoted \texttt{Rx}) that are automatically updated with a pre- or post-increment value \(k\). Address generation is based on modulo arithmetic. The programmer stores the length of the circular buffer into a \textit{modifier} (e.g., \texttt{Mx}) or \textit{modulo register} within the \gls{agu}. Each modifier is associated with at least one base register. While the starting address of the buffer can theoretically be arbitrary, circular buffers usually must be aligned to \(2^x\)-word boundaries, where \(2^x \geq \text{buffer size}\).

\paragraph{Special Addressing Mode: Bit-Reversed}
Bit-reversed addressing is specifically targeted to speed up \glsxtrfull{fft} algorithms. The output (or input, depending on the decimation strategy) of an \gls{fft} algorithm maps to bit-reversed array indices (e.g., index \(4_{10} = 100_2 \Longleftrightarrow 001_2 = 1_{10}\)). This mechanism is implemented in hardware via a real bit-reversal routing block or a reverse-carry arithmetic adder within the \gls{agu}.

\subsection{The TI-TMS320C54x \glsfmtshort{dsp}}
The Texas Instruments TMS320C54x is a classic 16-bit fixed-point \gls{dsp} designed specifically for low-power, high-performance telecommunications applications. Its architecture is built around an advanced modified Harvard bus structure, allowing concurrent instruction and multi-operand data fetches. The main characteristics of this processor include:
\begin{itemize}
    \item \textbf{Modified Harvard Architecture:} Features dual-access memory managed through multiple 16-bit buses (including a program bus and dedicated data buses like CB and DB), allowing parallel instruction fetching and dual data loads.
    \item \textbf{40-bit ALU:} While it is fundamentally a 16-bit processor, the \gls{alu} is 40 bits wide. This accommodates 32-bit results from the multiplier, while reserving 8 extra ``guard bits'' to prevent numerical overflow during prolonged \gls{mac} accumulation loops.
    \item \textbf{Hardware Multiplier:} A dedicated $16 \times 16$-bit hardware multiplier that generates a 32-bit product in a single clock cycle.
    \item \textbf{Dual 40-bit Accumulators:} Two independent accumulators allow parallel processing and temporary storage of wide mathematical results without having to route them back to main memory.
    \item \textbf{Advanced Multiplexing:} Most operand combinations can be fed directly to the \gls{alu} inputs through appropriate internal multiplexers, bypassing the register file and avoiding bottleneck delays.
\end{itemize}
\chapter{Pipelining}
This chapter explores the fundamental architectural concept of pipelining, a key implementation technique used to build high-performance processors.

\begin{definition}{Pipelining}{pipelining}
    Pipelining is an implementation technique where multiple instructions are overlapped in execution. It divides the elementary operations related to a single instruction into a sequence of smaller, discrete tasks (stages) that \textbf{can be performed in parallel} on different instructions simultaneously.
\end{definition}

\noindent This method allows the hardware execution units to process different elementary functions of different instructions concurrently. To maximize efficiency, the primary architectural goal is to balance the duration (length) of each pipeline stage. If the stages are perfectly balanced, the average time required per instruction when the pipeline is completely full approaches:
\[
    T_{\text{pipelined instruction}} = \frac{T_{\text{unpipelined instruction}}}{N_{\text{stages}}}
\]
This means that pipelining fundamentally decreases the average \gls{cpi} (ideally bringing it close to 1.0) without changing the underlying mathematical operations.

\begin{concept}{Ideal Pipeline Performance and Latency}{ideal-pipeline}
    Ideally, if the pipeline consists of \(p\) perfectly balanced stages and \(T_i\) is the total execution latency of an unpipelined instruction, then the delay of an individual pipeline stage (\(t_s\)) is:
    \[
        t_s = \frac{T_i}{p} + t_o
    \]
    where \(t_o\) is the overhead latency introduced by each stage boundary (e.g., pipeline register setup time, propagation delay, and clock skew). 

    \tcblower

    \textbf{Latency vs. Throughput:}\\
    It is important to note that pipelining \textbf{does not} reduce the total time it takes to execute a single instruction from start to finish; in fact, due to the overhead \(t_o\), the total latency slightly increases. 

    Instead, because the execution clocks are shared and overlapped, the effective \gls{cpi} is defined as the number of clock cycles required (on average) to complete an instruction \textit{after the previous one has completed}. Hence, once a pipeline is filled, one instruction finishes every \(t_s\) cycle. Consequently, the \textbf{instruction throughput} of a pipelined architecture is roughly \(p\) times larger than that of an equivalent unpipelined solution.
\end{concept}

\noindent From now forward, we will consider a theoretical \gls{risc} \gls{cpu} featuring single-sized instructions (32 bits wide). We assume the \gls{isa} consists exclusively of \gls{alu} operations, Load/Store operations, and Branches. Furthermore, \gls{alu} instructions operate exclusively on general-purpose registers (64 bits wide), and \textbf{only} Load/Store instructions are permitted to access main memory.

\section{Instruction Execution Stages}
Every instruction in this proposed architecture can be executed in at most five distinct clock cycles, mapping to the classic five-stage \gls{risc} pipeline:

\paragraph{Instruction Fetch (IF)} 
The \gls{pc} is sent to memory to fetch the current instruction from main memory (or the Instruction Cache). Simultaneously, the \gls{pc} is updated (incremented) to point to the next sequential instruction.

\paragraph{Instruction Decode / Register Fetch (ID)} 
The instruction is decoded while simultaneously reading the required operands from the register file. This simultaneous read is possible because \gls{risc} architectures utilize fixed-field instruction encoding (meaning the register specifiers are always located in the exact same bit positions). If the instruction utilizes an immediate value, this value is extracted and sign-extended.

\paragraph{Execution / Effective Address (EX)} 
The \gls{alu} operates on the operands prepared in the previous cycle, performing different tasks depending on the specific instruction class:
\begin{itemize}
    \item \textbf{Memory Reference (Load/Store):} The \gls{alu} adds the base register value and the sign-extended offset to calculate the effective memory address.
    \item \textbf{Register-Register (R-R) \gls{alu}:} The \gls{alu} performs the mathematical or logical operation specified by the opcode on the two values read from the register file.
    \item \textbf{Register-Immediate (R-I) \gls{alu}:} The \gls{alu} performs the operation between the register value and the sign-extended immediate value.
    \item \textbf{Branch:} The \gls{alu} adds the next instruction's address (\gls{pc}) to the sign-extended immediate offset to calculate the \glsxtrfull{bta}. Simultaneously, the branch condition is evaluated (often via a dedicated comparator).
\end{itemize}

\paragraph{Memory Access (MEM)} 
For Load instructions, data is read from the effective memory address calculated during the EX stage. For Store instructions, the data fetched from the register file is written into memory at the calculated effective address. For conditional Branch instructions, if the condition evaluated in the EX stage is met, the \gls{pc} is overwritten by the calculated \gls{bta}. \textit{(Note: If the instruction is a standard \gls{alu} operation, this stage remains idle).}

\paragraph{Write-Back (WB)} 
The final result is written back into the register file. For \gls{alu} instructions, this is the computed output from the EX stage; for Load instructions, this is the data fetched from memory during the MEM stage. \textit{(Store and Branch instructions do nothing during this stage).}

\begin{table}[!htbp]
    \centering
    \renewcommand{\arraystretch}{1.5} % Adds vertical breathing room between stages
    \begin{tabular}{| >{\raggedright\arraybackslash}p{0.08\textwidth} | >{\raggedright\arraybackslash}p{0.26\textwidth} | >{\raggedright\arraybackslash}p{0.26\textwidth} | >{\raggedright\arraybackslash}p{0.26\textwidth} |}
         \hline
         \textbf{Stage} & \textbf{ALU Instruction} & \textbf{Load/Store Instruction} & \textbf{Branch Instruction} \\
         \hline\hline
         
         \textbf{IF} & 
         \(\text{IR} \leftarrow \text{Mem}[\text{PC}]\) \newline \(\text{NPC} \leftarrow \text{PC} + 4\) & 
         \(\text{IR} \leftarrow \text{Mem}[\text{PC}]\) \newline \(\text{NPC} \leftarrow \text{PC} + 4\) & 
         \(\text{IR} \leftarrow \text{Mem}[\text{PC}]\) \newline \(\text{NPC} \leftarrow \text{PC} + 4\) \\
         \hline
         
         \textbf{ID} & 
         \(\text{A} \leftarrow \text{Rs1}; \quad \text{B} \leftarrow \text{Rs2}\) \newline \(\text{Imm} \leftarrow \text{sign\_extend}(\text{IR})\) & 
         \(\text{A} \leftarrow \text{Rs1}; \quad \text{B} \leftarrow \text{Rs2}\) \newline \(\text{Imm} \leftarrow \text{sign\_extend}(\text{IR})\) & 
         \(\text{A} \leftarrow \text{Rs1}; \quad \text{B} \leftarrow \text{Rs2}\) \newline \(\text{Imm} \leftarrow \text{sign\_extend}(\text{IR})\) \\
         \hline
         
         \textbf{EX} & 
         \(\text{ALUout} \leftarrow \text{A op B}\) \newline \textit{or} \(\text{ALUout} \leftarrow \text{A op Imm}\) & 
         \(\text{ALUout} \leftarrow \text{A} + \text{Imm}\) & 
         \(\text{ALUout} \leftarrow \text{NPC} + (\text{Imm} \times 4)\) \newline \(\text{Cond} \leftarrow (\text{A} == 0)\) \\
         \hline
         
         \textbf{MEM} & 
         \textit{(Idle)} & 
         \(\text{LMD} \leftarrow \text{Mem}[\text{ALUout}]\) \textit{(Load)} \newline \textit{or} \(\text{Mem}[\text{ALUout}] \leftarrow \text{B}\) \textit{(Store)} & 
         \(\text{if (Cond) PC} \leftarrow \text{ALUout}\) \\
         \hline
         
         \textbf{WB} & 
         \(\text{Rd} \leftarrow \text{ALUout}\) & 
         \(\text{Rd} \leftarrow \text{LMD}\) & 
         \textit{(Idle)} \\
         \hline
    \end{tabular}
    \caption{Register Transfer Language (RTL) definitions for the 5-Stage \gls{risc} Pipeline}
    \label{tab:pipeline-rtl}
\end{table}

\noindent Note that this baseline, unpipelined architecture executes each instruction in at most 5 clock cycles under the supervision of a simplified \gls{cu}. At the end of every clock cycle, all computed values are systematically written into an architectural or intermediate storage element (e.g., main memory, \glspl{gpr}, the \gls{pc}, or temporary pipeline registers). It is important to emphasize that this specific datapath is neither the only possible implementation of our \gls{isa}, nor the most efficient.

\paragraph{Unpipelined CPI Computation Example} 
Assume a processor operating with a \(1\,\text{ns}\) clock cycle, where \gls{alu} and Branch instructions take 4 cycles to complete, and Memory (Load/Store) instructions require the full 5 cycles. 

Given a typical instruction frequency distribution of:
\begin{itemize}
    \item \textbf{\gls{alu}:} $40\%$
    \item \textbf{Branch:} $20\%$
    \item \textbf{Memory:} $40\%$
\end{itemize}

\noindent We can calculate the average \gls{cpi} and the average instruction execution time (\(\overline{T_{\text{inst}}}\)) as follows:
\begin{align*}
    \overline{\text{CPI}} &= \left( (0.40 + 0.20) \times 4 \right) + (0.40 \times 5) \\
                          &= 2.4 + 2.0 = \mathbf{4.4 \text{ cycles/instruction}} \\[1em]
    \overline{T_{\text{inst}}} &= \text{Clock Period} \times \overline{\text{CPI}} \\
                               &= 1\,\text{ns} \times 4.4 = \mathbf{4.4\,\text{ns}}
\end{align*}

\section{Pipelined Execution and Speedup}
To visualize how pipelining improves instruction throughput, we use a space-time diagram (also known as a pipeline execution diagram). This maps the sequence of instructions (space) against the progression of clock cycles (time). 

In an ideal scenario without any interruptions, once the pipeline is full, every hardware stage is actively processing a different instruction simultaneously.

\begin{table}[!htbp]
    \centering
    \renewcommand{\arraystretch}{1.4}
    \begin{tabular}{|c|c|c|c|c|c|c|c|c|c|}
        \hline
        \textbf{Instruction} & \multicolumn{9}{c|}{\textbf{Clock Cycle Number}} \\
        \cline{2-10}
        \textbf{Order} & \textbf{1} & \textbf{2} & \textbf{3} & \textbf{4} & \textbf{5} & \textbf{6} & \textbf{7} & \textbf{8} & \textbf{9} \\
        \hline\hline
        \(i\)     & IF & ID & EX & MEM & WB &    &    &    &    \\
        \hline
        \(i+1\)   &    & IF & ID & EX & MEM & WB &    &    &    \\
        \hline
        \(i+2\)   &    &    & IF & ID & EX & MEM & WB &    &    \\
        \hline
        \(i+3\)   &    &    &    & IF & ID & EX & MEM & WB &    \\
        \hline
        \(i+4\)   &    &    &    &    & IF & ID & EX & MEM & WB \\
        \hline
    \end{tabular}
    \caption{Basic 5-Stage Pipelined \gls{risc} Space-Time Diagram (Ideal Execution)}
    \label{tab:pipeline-diagram}
\end{table}

\noindent Notice that during Clock Cycle 5, the pipeline is fully saturated: the \gls{cpu} is simultaneously writing back instruction \(i\), accessing memory for \(i+1\), executing \(i+2\), decoding \(i+3\), and fetching \(i+4\).

\paragraph{Pipelined CPI and Speedup Computation}
In this ideal pipeline, there are no stalls (no structural, data, or control hazards). After the initial \(p-1\) cycles required to fill the pipeline (where \(p=5\) stages), exactly one instruction finishes execution during every single clock cycle. Thus, the ideal pipelined \gls{cpi} remains strictly $1.0$.

However, physical pipelining requires the insertion of intermediate pipeline registers between stages. These registers introduce unavoidable hardware overhead (\(t_o\)), such as register setup time, propagation delay, and clock skew. Assuming this overhead forces our clock cycle period to increase from the unpipelined \(1.0\,\text{ns}\) to \(1.2\,\text{ns}\), the pipelined performance is calculated as follows:
\begin{align*}
    \overline{\text{CPI}}_{\text{pipelined}} &= 1.0 \text{ cycle/instruction} \\
    \overline{T_{\text{pipelined}}} &= \text{Clock Period} \times \overline{\text{CPI}}_{\text{pipelined}} \\
                                    &= 1.2\,\text{ns} \times 1.0 = \mathbf{1.2\,\text{ns}}
\end{align*}

\noindent Using the unpipelined architecture from the previous section as our baseline (\(\overline{T_{\text{unpipelined}}} = 4.4\,\text{ns}\)), we can calculate the realistic speedup achieved. Note that while the theoretical maximum speedup is \(5\times\) (the number of stages), our actual speedup is lower due to both the unpipelined \gls{cpu} already optimizing some instructions to 4 cycles, and the hardware overhead added to the clock period:
\[
    \text{Speedup} = \frac{\overline{T_{\text{unpipelined}}}}{\overline{T_{\text{pipelined}}}} = \frac{4.4\,\text{ns}}{1.2\,\text{ns}} \approx \mathbf{3.67\times}
\]

\noindent This calculation beautifully illustrates the core trade-off of \gls{risc} pipelining: we accept hardware overhead (a \(1.2\,\text{ns}\) clock instead of \(1.0\,\text{ns}\)) and force faster instructions to pass through idle stages so that the datapath remains perfectly synchronized. In return, we achieve a massive overall throughput speedup of \(3.67\times\).

\section{Hazards and Stalls}
The presence of several instructions simultaneously active within the pipeline may lead to various types of dependencies which turn into \textbf{hazards}. These reduce overall performance by forcing the pipeline to \textbf{stall} (insert idle cycles, often called ``bubbles'') until the source of the hazard is resolved. Hazards are fundamentally divided into three categories:
\begin{description}
    \item[Structural:] Arise from hardware resource conflicts when multiple instructions require the same functional unit simultaneously.
    \item[Data:] Occur when an instruction depends on the computed result of a previous, still-in-flight instruction.
    \item[Control:] Arise from the pipelining of branches or other instructions that modify the \gls{pc}.
\end{description}

\subsection{Structural Hazards}
A structural hazard arises when two instructions attempt to use the same hardware component at the same time (e.g., if one instruction wants to write to memory in its \texttt{MEM} stage while another wants to fetch an instruction from memory in its \texttt{IF} stage). 

To solve this, the simplest solution is to stall the subsequent instruction until the hardware hazard has passed. However, this inevitably produces a reduction in throughput. More advanced solutions involve hardware resource splitting/replication (such as splitting a unified memory into separate Instruction and Data caches) or software-level instruction reordering by the compiler. 

The negative impact on the \gls{cpi} is strictly proportional to the probability of a stall. Assuming a 20\% stall rate (i.e., $1/5$), it is calculated as follows:
\[
    \text{\gls{cpi}} = \frac{N_{\text{clocks}}}{N_{\text{inst}}} = \frac{\gls{ic-inst} + \left(\frac{\gls{ic-inst}}{5} - 1\right)}{\gls{ic-inst}} = \frac{6}{5} - \frac{1}{\gls{ic-inst}}
\]
For an asymptotically large \gls{ic-inst}, this limits the performance to:
\[
    \lim_{\gls{ic-inst} \to \infty} \text{\gls{cpi}} = 1.2
\]

\subsection{Data Dependencies and Hazards}
Data hazards are strictly related to data dependencies. Two instructions are data-dependent if they access a common register or memory location (unless both instructions are exclusively reading the data). Data dependencies can occur in straight-line code (subsequent instructions) or across loop iterations. There are three kinds of data dependencies:
\begin{itemize}
    \item \textbf{\glsxtrfull{raw}:} Flow or true dependence.
    \item \textbf{\glsxtrfull{war}:} Anti-dependence.
    \item \textbf{\glsxtrfull{waw}:} Output dependence.
\end{itemize}

While \gls{war} and \gls{waw} are considered \textit{false dependencies} (which can often be resolved in modern architectures through hardware register renaming), \gls{raw} hazards are \textit{true dependencies}. They cannot be eliminated, but they can be mitigated in certain situations by instruction scheduling or dedicated hardware routing. To avoid data hazards, the \gls{cu} must dynamically:
\begin{enumerate}
    \item Check the source registers (\(Rs_1, Rs_2\)) of the current instruction against the destination register (\(Rd\)) of all older, incomplete instructions currently in the pipeline (requiring up to 6 simultaneous hardware comparisons).
    \item If a hazard is detected, the \gls{cu} must stall the conflicting instruction in the \texttt{ID} stage until the required data is written back.
\end{enumerate}

The impact of unmodified data hazards on the \gls{cpi} can be tremendous. Assuming one data dependency occurs every six operations, and each causes three stall cycles, the performance degrades significantly:
\[
    \text{\gls{cpi}} = 1 + \frac{3}{6} = 1.5
\]
This represents a massive 50\% performance penalty compared to the ideal pipeline \gls{cpi} of 1.0.

\paragraph{Stall Penalty Reduction: Forwarding (Bypassing)} 
This stall penalty can be drastically reduced through a hardware technique called \textbf{forwarding} (or bypassing). Because source operands are typically only necessary at the beginning of the \texttt{EX} stage for instruction \(I_{k+1}\), and the required destination result is already computed by the end of the \texttt{EX} stage for instruction \(I_k\), the output can be re-routed. A hardware bypass routes the result directly back to the \gls{alu} input, completely skipping the delay of waiting for the \texttt{WB} stage.

To implement this, the source inputs of the \gls{alu} must be managed by multiplexers (MUXes). The \gls{cu} switches these MUX inputs from the standard \texttt{ID/EX} pipeline registers to either the \texttt{EX/MEM} or \texttt{MEM/WB} pipeline registers, dynamically routing the freshest data based on the proximity of the dependent instructions.

\subsection{Control Dependencies and Hazards}
Similar to data hazards, control hazards are strictly related to control dependencies. These are sequential constraints dictated by the control flow of the program. All branch instructions create a control dependency, and the ratio between branch instructions and straight-line code heavily depends on the specific application being executed. 

\noindent If a conditional branch changes the \gls{pc} to the target address, it is referred to as a \textbf{taken} branch; if it falls through to the next sequential instruction, it is \textbf{not taken} (or untaken). Control hazards arise from the risk of fetching the wrong instructions into the pipeline following a conditional branch before the branch condition is actually resolved. These hazards can seriously decrease \gls{cpu} performance.

\paragraph{Simple Control Hazard Solutions} 
There are several methods to mitigate control hazards. The simplest four approaches are:
\begin{itemize}
    \item \textbf{Stall / Freeze:} Hold or delete any instruction fetched after the branch until the \gls{bta} and branch outcome are completely resolved.
    \item \textbf{Predict Not Taken:} Assume the branch will fall through. If it is actually taken, flush the incorrectly fetched instructions.
    \item \textbf{Predict Taken:} Assume the branch will be taken and fetch from the target (requires knowing the \gls{bta} immediately). If it is not taken, flush the instructions.
    \item \textbf{Delayed Branch:} The compiler reorders the code to move independent instructions (which must be executed regardless of the branch outcome) into the ``delay slots'' immediately following the branch instruction.
\end{itemize}

Although theoretical branch delay slots can be multiple instructions long, in practice, it is difficult for compilers to find enough independent instructions to fill more than one delay slot. To reduce the control hazard time penalty entirely, architects can shift the branch condition evaluation and the \gls{bta} calculation from the \texttt{EX} stage back to the \texttt{ID} stage. This requires adding a dedicated adder and comparator to the \texttt{ID} stage hardware, but it successfully decreases the branch penalty from three stalls down to just one.

\subsection{\glsfmtshort{cpi} and Speedup after Hazards}
\begin{concept}{Realistic CPI and Speedup with Hazards}{realistic-cpi}
    The actual speedup of a pipelined processor, explicitly accounting for the delays introduced by data and control hazards, is calculated as follows:
    \[\text{Speedup} = \frac{\overline{\text{CPI}}_{\text{unpipelined}}}{\overline{\text{CPI}}_{\text{pipelined}}} \times \frac{T_{\text{clk, unpipelined}}}{T_{\text{clk, pipelined}}}\]
    Assuming \(\overline{N}_{\text{stalls}}\) represents the average number of stall cycles introduced per instruction due to hazards, the pipelined \gls{cpi} is:
    \[\overline{\text{CPI}}_{\text{pipelined}} = \text{Ideal CPI} + \overline{N}_{\text{stalls}} = 1.0 + \overline{N}_{\text{stalls}}\]
    If we ignore the clock cycle hardware overhead (assuming \(T_{\text{clk, unpipelined}} = T_{\text{clk, pipelined}}\)) and assume all instructions in the unpipelined baseline architecture take a number of cycles equal to the pipeline depth, the simplified speedup equation becomes:
    \[\text{Speedup} = \frac{\overline{\text{CPI}}_{\text{unpipelined}}}{1 + \overline{N}_{\text{stalls}}} = \frac{\text{Pipeline Depth}}{1 + \overline{N}_{\text{stalls}}}\]
\end{concept}

\section{Exceptions}
The term ``exceptions'' is a broad category that covers all exceptional situations whereby the normal execution order of instructions is disrupted. The management of these events varies significantly among different \glspl{cpu}. 

There are numerous sources of exceptions, including:
\begin{itemize}
    \item Conditions arising from instruction execution (e.g., arithmetic overflow, floating-point anomalies).
    \item External interrupts (e.g., I/O requests).
    \item System calls (traps).
    \item Memory-related exceptions (e.g., page faults, protection violations).
    \item Undefined opcodes.
    \item Hardware malfunctions and power failures.
\end{itemize}

\noindent Not all \glspl{cpu} are designed to support every exception type natively; some higher-level exceptions, such as breakpoints or instruction tracing, are exclusively implemented for debugging purposes.

\paragraph{Exception Handling in Non-Pipelined \glspl{cpu}} 
In a non-pipelined architecture, execution management is relatively simple. Because only one instruction executes at a time, an exception is unambiguously associated with that specific instruction. These events are managed by dedicated routines. When an exception arises, the \gls{cpu} looks up a vector table containing the memory address of the corresponding Exception Service Routine (ESR). The hardware saves the current Program Counter (\gls{pc}), runs the ESR to handle the exception, and then returns control to the original program if required.

\paragraph{Exception Handling in Pipelined \glspl{cpu}} 
In pipelined architectures, the situation is drastically more complicated. The processor must correctly identify exactly which instruction in the pipeline caused the exception and save the precise machine state at that exact moment. Because different types of exceptions occur at different stages of the pipeline, the management unit must utilize distinct handling procedures. The most difficult scenario involves handling exceptions where the faulting instruction must be resumed after the exception is resolved (e.g., a page fault); this requires a completely \textbf{restartable pipeline}.

\paragraph{Pipeline Stages and Specific Exceptions}
Exceptions are generally tied to the specific pipeline stage in which they are detected:
\begin{itemize}
    \item \textbf{\texttt{ID} Stage (Undefined Opcode):} This may be due to a compiler error, a memory fault, or invoking functional units not present in the specific device model. This exception is non-maskable and usually leads to program termination.
    \item \textbf{\texttt{EX} Stage (Execution Anomalies):} Arithmetic exceptions, such as overflow, are raised here. They are generally user-maskable and require the program to be resumed.
    \item \textbf{\texttt{MEM} Stage (Memory Faults):} Data cache misses and page faults are detected here. These exceptions are non-maskable and require a restartable pipeline.
    \item \textbf{\texttt{WB} Stage / Asynchronous (External Interrupts):} I/O interrupts are asynchronous but require resumption. The Interrupt Service Routine (ISR) is usually serviced when the interrupted instruction reaches the \texttt{WB} stage.
\end{itemize}

\paragraph{Simultaneous and Out-of-Order Exceptions}
Because multiple instructions are in flight simultaneously, exceptions may occur concurrently or out of program order:
\begin{itemize}
    \item \textbf{Simultaneous:} A \texttt{LD} instruction raises a page fault in the \texttt{MEM} stage at the exact same clock cycle that a subsequent \texttt{DADD} instruction raises an arithmetic exception in the \texttt{EX} stage.
    \item \textbf{Out-of-Order:} A \texttt{LD} instruction is in the \texttt{MEM} stage, and a subsequent \texttt{DADD} instruction is in the \texttt{IF} stage. If \texttt{DADD} raises a page fault (instruction fetch), its exception actually occurs \textit{in time} before the \texttt{LD} raises a data page fault, even though \texttt{LD} precedes \texttt{DADD} in the program order.
\end{itemize}

\paragraph{The Pipeline Exception Handling Protocol}
To solve the complexities of out-of-order and simultaneous exceptions, pipelined \glspl{cpu} do not handle exceptions immediately. Instead, they post the exception in an \textbf{Exception Status Vector} that travels alongside the instruction down the pipeline. 

The strict protocol for handling exceptions is as follows:
\begin{enumerate}
    \item \textbf{Halt State Modification:} Turn off both register and memory writes for the faulting instruction and \textit{all instructions following it} in the pipeline. All instructions preceding the faulting one are allowed to complete normally and modify the machine state (Registers, Memory). All subsequent instructions are flushed unexecuted.
    \item \textbf{Check Status at \texttt{WB}:} When an instruction enters the \texttt{WB} stage, the hardware checks its exception status vector. If any exceptions are posted, they are finally handled. This guarantees that exceptions are serviced in strict program order, regardless of when they were physically detected.
    \item \textbf{Save and Transfer Control:} After the exception-handling routine receives control, it saves the \gls{pc} of the faulting instruction so the processor can accurately return from the exception later on.
\end{enumerate}

\section{Advanced pipelining}
Beyond simply shortening the clock cycle, modern \glspl{cpu} improve performance by increasing hardware parallelism. This increases the ideal number of instructions executed per cycle and reduces structural hazards. Modern processors achieve this by combining two complementary approaches:

\begin{itemize}
    \item \textbf{Superpipelining:} The logical pipeline steps (\texttt{IF, ID, EX, MEM, WB}) are subdivided into even smaller, simpler stages, allowing for a much shorter clock cycle ($T_{CLK}$).
    \item \textbf{Superscalar Architecture:} The processor features multiple identical or dedicated physical pipelines, allowing it to decode, issue, and execute multiple scalar operations in parallel in a single clock cycle.
\end{itemize}

\paragraph{The Limits of Superpipelining}
While superpipelining reduces the clock cycle time, it increases the total number of cycles required to complete a program. Furthermore, the depth of the pipeline (usually capped at 12--15 stages in modern architectures) is strictly limited by four factors:
\begin{enumerate}
    \item \textbf{Imbalance Problems:} The clock period can never be shorter than the slowest sub-stage.
    \item \textbf{Pipeline Overhead:} The physical delays of the hardware (register setup/propagation time and clock skew) create a hard lower bound for the clock cycle.
    \item \textbf{Growing Data Dependences:} A deeper pipeline means data hazards result in a higher number of inserted bubbles, increasing the performance penalty.
    \item \textbf{Longer Control Hazards:} Wrong branch predictions are heavily penalized because a much larger number of stages must be flushed.
\end{enumerate}

\paragraph{Handling Multicycle Operations}
Certain operations, such as floating-point arithmetic, cannot physically complete in a single cycle. To handle this, the \texttt{EX} stage is repeated as many times as necessary, and the \gls{cpu} employs multiple dedicated functional units (e.g., FP Add, FP Multiply, FP Divide). 

However, allowing multicycle operations means instructions no longer finish in strict program order, altering how hazards behave:
\begin{itemize}
    \item \textbf{WAW (Write-After-Write) Hazards:} These are now possible because a faster instruction might reach the \texttt{WB} stage before an older, slower multicycle instruction.
    \item \textbf{WAR (Write-After-Read) Hazards:} These remain impossible because all operands are still strictly read in order during the \texttt{ID} stage.
    \item \textbf{RAW (Read-After-Write) Stalls:} These become much more frequent due to the longer latency of the multicycle instructions.
    \item \textbf{Structural Hazards:} Stalls will occur if multiple instructions attempt to use the same functional unit, or if they arrive at the \texttt{MEM} or \texttt{WB} stages simultaneously.
\end{itemize}

\paragraph{Complexities of Superscalar Execution}
To extract true parallelism from a program, superscalar architectures require advanced hardware management. Because multiple pipelines are executing instructions of varying lengths simultaneously, the \gls{cpu} relies heavily on \textbf{dynamic scheduling}, \textbf{branch prediction}, and \textbf{speculative execution}.

This out-of-order execution environment makes maintaining precise exceptions incredibly difficult. Consequently, these architectures may suffer from \textbf{imprecise exceptions}—where the saved state of the processor at the time of an exception does not perfectly match the strict sequential order of the original code.
\chapter{Memory Organization}

Applications require, year after year, larger memory capacities and higher overall performance. However, constructing a single memory unit that is simultaneously large, fast, and inexpensive is physically and economically impossible: larger memories are inherently slower, and faster memories are significantly more expensive. 

Furthermore, historical trends show a widening performance gap. On average, \gls{dram} latency has decreased by only about 7\% per year, whereas \gls{cpu} performance has grown by up to 135\% per year. This disparity means the performance gap between processors and main memory has grown significantly over time, creating a severe bottleneck.

\paragraph{Solution: Memory Hierarchy}
To overcome this limitation, modern computing systems organize memory into a \textbf{hierarchy}. This hierarchy is composed of multiple levels characterized by different manufacturing technologies, costs, physical dimensions, and access mechanisms. 

\begin{itemize}
    \item \textbf{Higher levels} (closer to the \gls{cpu}) are smaller, faster, and more expensive per byte.
    \item \textbf{Lower levels} (further from the \gls{cpu}) are much larger and cheaper, but significantly slower.
\end{itemize}

This architectural abstraction provides the best of both worlds: the programmer perceives a single, vast memory space, while the \gls{cpu} interacts with a fast memory system built at an acceptable overall cost. To function correctly, this hierarchy must implement efficient, run-time data transfer mechanisms between its levels. These mechanisms rely heavily on a fundamental behavioral trait of computer programs known as the \textit{Principle of Locality}.

\begin{concept}{The Principle of Locality}{locality}
    Programs do not access memory addresses randomly. Instead, they exhibit specific patterns of access, which memory hierarchies exploit to ensure data is moved to the fastest levels just before the \gls{cpu} needs it. Locality is strictly divided into two typologies:
    \begin{itemize}
        \item \textbf{Temporal Locality:} If a memory location is referenced recently, there is a high probability it will be referenced again in the near future (e.g., loop variables).
        \item \textbf{Spatial Locality:} If a memory location is referenced, there is a high probability that neighboring memory locations will be referenced shortly after (e.g., array iteration).
    \end{itemize}
\end{concept}

\noindent Taking these principles into account, memory hierarchies are explicitly designed to exploit both temporal and spatial locality to ensure the quickest possible memory access.

\paragraph{Locality Variations in Code}
The manifestation of locality can vary depending on whether the processor is fetching instructions or reading/writing data:

\subparagraph{Instruction Locality} 
This is generally highly predictable and application-independent. Because code normally executes sequentially and relies heavily on loops, it exhibits excellent spatial and temporal locality. A common metric is the \textbf{90/10 Rule}, which states that a program typically spends 90\% of its execution time running only 10\% of its actual code.

\subparagraph{Data Locality} 
This is highly application-dependent. While some applications might exhibit strong data locality, others might show poor data locality overall, or heavily favor one typology over the other (e.g., streaming large media files exhibits high spatial locality but zero temporal locality).

\paragraph{Types of Memory Hierarchies} 
Different computing needs dictate entirely different memory hierarchy priorities and features.

\subparagraph{Desktop Systems} 
These systems typically run one or more applications on top of the \gls{os} for a single user. The primary concern here is minimizing the \textbf{average latency} of the memory hierarchy to ensure a responsive user experience. They still require traditional mass storage at the lowest tier for permanent data retention.

\subparagraph{Server Systems} 
Servers usually handle hundreds of concurrent users and run many heavy applications simultaneously. Consequently, the focus shifts toward minimizing the cost of context switching rather than strictly minimizing average latency. Servers also heavily prioritize overall memory bandwidth, data protection (e.g., ECC memory), and fault tolerance.

\subparagraph{Embedded Systems} 
These solutions require a completely different approach based on their specific physical constraints and use cases:
\begin{itemize}
    \item \textbf{Real-Time Applications:} Must prioritize guaranteed \textbf{worst-case latency} over average latency to ensure critical deadlines are always met.
    \item \textbf{Mobile/IoT Devices:} Are strictly constrained by power consumption and battery life.
\end{itemize}
Embedded architectures usually run a single application alongside a very small, lightweight \gls{os} (or no \gls{os} at all). For these reasons, their main memory is much smaller, and traditional spinning disk storage is largely unnecessary.

\vspace{1em}
\begin{table}[H]
    \centering
    \renewcommand{\arraystretch}{1.4}
    \resizebox{\textwidth}{!}{%
    \begin{tabular}{r | l l l l}
        \textbf{Characteristics} & \textbf{Registers} & \textbf{Cache} & \textbf{Main Memory} & \textbf{Disk Storage} \\
        \hline
        \textbf{Size} & $< 1$ KB & $< 16$ MB & $< 512$ GB & $> 1$ TB \\
        \textbf{Technology} & Custom multi-port memory & On/Off-chip CMOS SRAM & CMOS DRAM & Magnetic Disk / SSD \\
        \textbf{First Access Time (ns)} & 0.25 -- 0.5 & 0.5 -- 25 & 50 -- 250 & $\sim 5,000,000$ \\
        \textbf{Bandwidth (MB/s)} & 50,000 -- 500,000 & 5,000 -- 20,000 & 2,500 -- 10,000 & 50 -- 500 \\
        \textbf{Managed By} & Compiler & Hardware & \gls{os} & \gls{os} / Operator \\
        \textbf{Backed By} & Cache & Main Memory & Disk & CD / Tape (Archive) \\
    \end{tabular}%
    }
    \caption{Performance and implementation characteristics across the memory hierarchy.}
\end{table}

\paragraph{The Register File} 
Registers are typically exposed directly to the programmer or managed strictly through the compiler. Their contents are stored in a centralized register file that features multiple read and write ports to facilitate simultaneous parallel data access. 

A standard three-operand instruction requires a register file with two read ports and one write port. However, as the number of parallel functional units within a processor increases, the number of memory ports—and consequently the total number of registers—must scale accordingly. For example, if a \gls{cpu} features $N$ \glsxtrfullpl{alu}, its register file must theoretically support $2N$ read ports and $N$ write ports. 

Unfortunately, the physical silicon area of a register file grows quadratically with the square of the number of ports. This means an excessive number of ports will severely slow down the processor.

\subparagraph{Multi-Port Advantages and Limitations} 
This quadratic increase in chip area is primarily driven by the massive amount of additional input and output wiring required. Furthermore, the combination of complex signal routing and increased internal bit-cell loading results in higher power consumption and a linear increase in access delay proportional to the number of ports. 

Due to these severe hardware limitations, a register file is typically capped at a maximum of 15 to 20 ports. When an architecture employs more than 4 functional units, architectural trade-offs dictate that 12 ports offers the absolute optimal balance between parallel throughput, physical space, and access latency.

\section{On-Chip Cache and RAM Memories}
The first level of memory, typically located directly on-chip, is the fastest type of memory available (excluding the \gls{cpu} registers). This memory is constructed using high-speed \gls{sram} and serves as the primary storage for immediately required data and instructions. Due to its limited physical size, architects typically implement one of two primary memory solutions at this level: standard \gls{sram} or a hardware Cache.

\paragraph{Standard \gls{sram}} 
This type of memory is explicitly exposed to the programmer and the compiler. Data transfers must be performed manually by the software. Its use is generally restricted to low-end and specialized embedded processors, where strict worst-case latency guarantees and low power consumption are the primary concerns (often utilizing standard Von Neumann or Harvard architectures).

\paragraph{Scratch Pad Memories}
As a low-power alternative to complex caches, a scratch pad is a simple \gls{sram} memory buffer where the most frequently used program functions, instructions, and data are explicitly allocated. Unlike a cache, it lacks hardware controller support. Instead, it requires preliminary software profiling to determine exactly which instructions should reside in the pad. This allocation is performed statically at compile time. Scratch pads are widely used in embedded systems because these devices typically run a fixed program, making static allocation highly predictable and offering an excellent trade-off between performance and power consumption.

\paragraph{Hardware Cache} 
In contrast to standard \gls{sram} and scratch pads, a cache is completely transparent to software developers and compilers. Data transfers between the cache and lower memory levels are managed entirely by hardware (the internal control logic of the architecture itself). This solution is heavily favored in general-purpose computing, where optimizing the \textit{average} performance is the primary objective.

\subsection{Cache Basics}
A cache is a specialized memory array built from \gls{sram} cells that stores the instructions and data with the highest mathematical probability of being used in the near future (leveraging the Principle of Locality). 

When the \gls{cpu} requests a specific memory location, it checks the cache first. The cache control logic determines if the data is present. If it is not, the hardware controller automatically halts the \gls{cpu} to fetch the missing block from a lower memory level. These lower levels might be subsequent, progressively larger caches organized in a sub-hierarchy (e.g., L1, L2, \dots, LN) or the main \gls{dram}.

To evaluate the efficiency of a cache hierarchy, computer architects rely on the following standard metrics:

\begin{definition}{Cache Hit}{cache-hit}
    A memory access where the data or instruction requested by the \gls{cpu} is successfully found in the cache.
\end{definition}

\begin{definition}{Cache Miss}{cache-miss}
    A memory access where the requested item is not present in the cache, requiring a fetch from a lower memory level.
\end{definition}

\begin{definition}{Hit Rate}{hit-rate}
    The fraction of total memory accesses that successfully result in a cache hit.
\end{definition}

\begin{definition}{Miss Rate}{miss-rate}
    The fraction of memory accesses that result in a cache miss, mathematically defined as $1 - \text{Hit Rate}$.
\end{definition}

\begin{definition}{Hit Time}{hit-time}
    The time (or number of clock cycles) required to access the cache on a successful hit. This intrinsically includes the processing time needed to compare memory addresses and determine whether the access is a hit or a miss.
\end{definition}

\begin{definition}{Miss Penalty}{miss-penalty}
    The extra time (in clock cycles) required to fetch a missing block of data from lower-level storage and allocate it into the cache.
\end{definition}

\begin{definition}{Miss Time}{miss-time}
    The total overall time required to obtain the requested data in the event of a miss, calculated as $\text{Hit Time} + \text{Miss Penalty}$.
\end{definition}

\subsection{Cache Structure and Design Problems}
At a hardware level, a cache is structured as an array of \textbf{blocks} (sometimes called lines), where each block contains a fixed number of words from memory. To simplify addressing and hardware decoding, the size of a block is always a power of 2 (e.g., 32 words). 

When the \gls{cpu} requests a specific word, its entire enclosing block is loaded from lower-level memory and placed into the cache. Note that blocks sitting next to each other in main memory are not necessarily stored next to each other in the cache, and vice versa.

To determine whether a requested word is currently present, the cache cannot just check the address; it must store each block alongside its actual memory address tag. The hardware compares the requested address against the tags of the stored blocks simultaneously. If a matching tag is found and the data is valid, it is a \textbf{cache hit}; otherwise, it is a \textbf{cache miss}.

\paragraph{The Four Core Cache Design Problems}
Designing a hardware cache requires solving four fundamental engineering problems:
\begin{enumerate}
    \item \textbf{Placement Problem:} Where can a block be placed when it is transferred from a lower-level memory to a higher-level cache?
    \item \textbf{Search (Identification) Problem:} How does the hardware quickly determine whether a requested item is present in the cache?
    \item \textbf{Substitution (Replacement) Problem:} When a cache miss occurs and the cache is full, which existing block must be evicted to make room for the new block?
    \item \textbf{Write Strategy Problem:} How should the system handle write-to-memory instructions? \emph{(Note: This problem is orthogonal to the first three and will be treated separately once basic read architectures are established).}
\end{enumerate}

You are completely right. Cramming Direct-Mapped Caches into a single bullet point stripped away the precise mathematical definitions and clear prose that made it so readable. Direct-mapped caches are the foundation of understanding all other organizations, so they deserve their own dedicated subsection with full breathing room.

Here is the ideal balance: we restore your favorite, detailed breakdown of Direct-Mapped Caches in full, and then give N-Way Set-Associative and Fully Associative caches their own distinct, uncompacted subsections so the entire section maintains equal depth and clarity.
Snippet di codice

\subsection{Direct-Mapped Caches}
The simplest approach to the placement problem is the \textbf{direct-mapped cache}. 

Let $NB$ be the total number of blocks in the cache. Let $BAR$ (Block Address in RAM) be the main memory address stripped of its lower bits (i.e., removing the bits identifying the specific byte within the word and the word within the block). 

In a direct-mapped cache, each block from RAM can be placed in \textit{one and only one} specific location in the cache, calculated via modulo arithmetic:
\[ BAC = BAR \pmod{NB} \]
where $BAC$ is the Block Address in the Cache, commonly referred to as the \textbf{cache index}. If $NB$ is a power of 2 (i.e., $NB = 2^{NAC}$), then the cache index corresponds precisely to the $NAC = \log_2(NB)$ least significant bits of the $BAR$.

\paragraph{Solving Search and Placement in Direct-Mapped Caches}
Because of the strict modulo mapping, the placement problem is trivial: blocks have a fixed, predetermined destination. However, the search problem is more challenging because multiple different memory blocks from different parts of RAM map to the exact same cache slot (creating potential ``cache collisions'').

To solve this, each cache line must be equipped with two hardware fields:
\begin{itemize}
    \item \textbf{Tag Field:} Stores the remaining upper bits of the $BAR$ (specifically, the total address bits minus the index bits $NAC$). 
    \item \textbf{Validity Bit:} A single bit indicating whether the cache line currently contains meaningful, initialized data or is empty/invalid upon startup.
\end{itemize}
To verify if data is present, the hardware extracts the index bits of the requested address, looks up that specific cache line, checks if the validity bit is set, and compares the line's stored \textbf{Tag} with the upper bits of the requested address. If they match, it is a hit.

\paragraph{The Replacement Problem and Limitations}
The replacement problem is trivial in a direct-mapped cache: when a new block arrives, it has only one possible slot it can occupy, so it automatically evicts whatever block is currently residing there, with no hardware decision-making required. 

However, this introduces a major architectural disadvantage: \textbf{it completely ignores temporal locality!} If two frequently accessed variables happen to map to the same $BAC$ index, they will repeatedly evict each other on every access, causing a spike in conflict misses even if the rest of the cache is completely empty.

\subsection{N-Way Set-Associative Caches}
To bridge the gap between the strict hardware simplicity of direct-mapped caches and the heavy complexity of fully associative designs, modern systems heavily rely on \textbf{N-way set-associative caches}. 

The cache is organized into $N_{set}$ sets (where $N_{set}$ is a power of 2), with each set containing exactly $N$ blocks. Notice that the total number of blocks is $NB = N_{set} \times N$.
\begin{itemize}
    \item \textbf{Placement Problem:} The target set index ($I_{set}$) corresponding to a given memory block address $BAR$ is determined using modulo arithmetic:
    \[ I_{set} = BAR \pmod{N_{set}} \]
    Once the set is identified via its index, the block can be placed into \emph{any} of the $N$ cache blocks within that specific set using fully associative rules.
    \item \textbf{Search Problem:} The set itself is located just like a direct-mapped cache (using the index bits). Within that set, the hardware executes parallel tag comparisons across the $N$ blocks to find a match.
\end{itemize}

\paragraph{Architectural Trade-offs and Scaling}
\begin{itemize}
    \item A standard direct-mapped cache is simply a 1-way set-associative cache ($N = 1$).
    \item \textbf{Compared to Fully Associative Caches:} Set-associativity requires significantly fewer and smaller comparators, resulting in cheaper, faster, and more power-efficient hardware.
    \item \textbf{Compared to Direct-Mapped Caches:} It drastically reduces conflict misses by allowing multiple blocks to share the same set index, better exploiting temporal locality at a reasonable cost.
\end{itemize}

\noindent \textbf{Scaling Associativity (Constant Cache Size):}
If you double the degree of associativity ($N$) while keeping total cache size constant:
\begin{itemize}
    \item The number of blocks within a single set doubles $\implies$ the tag length grows by 1 bit, and the number of comparators per set doubles.
    \item The total number of sets ($N_{set}$) is halved $\implies$ the index length shrinks by 1 bit.
    \item \emph{Empirical Note:} While higher associativity reduces conflict misses, the performance gains diminish rapidly for large caches, making 4-way to 16-way designs the sweet spot for modern processors.
\end{itemize}

\subsection{Fully Associative Caches}
At the opposite extreme of the design spectrum is the \textbf{fully associative cache}, where every block of RAM can be mapped onto \emph{any} block line in the cache with zero placement constraints.
\begin{itemize}
    \item \textbf{Placement Problem:} Non-existent. The cache index size is $0$ bits, and the tag comprises the complete address minus the block and byte displacements.
    \item \textbf{Search Problem:} Because a block can reside anywhere, the hardware cannot use an index lookup. Instead, it must run a \textbf{parallel comparison} of the requested address against all tags in the cache simultaneously. This requires a massive array of comparators, making it extremely expensive in silicon area and power consumption.
    \item \textbf{Replacement Problem:} Because placement is completely flexible, the cache requires explicit replacement policies when a miss occurs and the cache is full:
    \begin{itemize}
        \item \textbf{Random:} The block to be substituted is chosen randomly. It is simple and requires minimal hardware (a pseudo-random number generator), but its miss-rate efficiency is poor because it ignores temporal locality.
        \item \textbf{Least Recently Used (LRU):} Evicts the block left unused for the longest time. It fully exploits temporal locality for maximum efficiency, but requires heavy hardware counter tracking whose complexity scales poorly with the number of blocks.
        \item \textbf{First-In, First-Out (FIFO) / Round-Robin:} Evicts the oldest block by entry order. It provides the best practical trade-off between performance and hardware complexity and is widely used in modern systems, though it only approximates locality.
    \end{itemize}
\end{itemize}

\subsection{Causes of Cache Misses}
To understand how cache organizations improve performance, misses are classified into three fundamental categories:
\begin{itemize}
    \item \textbf{Compulsory Misses:} The very first access to a block (cold-start misses). These are unavoidable because the data is not yet in the cache.
    \item \textbf{Capacity Misses:} Occur because the cache is physically too small to hold all the blocks needed during execution. Blocks are evicted and later re-fetched.
    \item \textbf{Conflict Misses (Collision Misses):} Occur in direct-mapped or set-associative caches when multiple blocks compete for the exact same cache line or set, causing premature eviction despite available space elsewhere.
\end{itemize}

\subsection{The Write Problem and Write Strategies}
As noted earlier, handling write instructions is orthogonal to read placement strategies. When the \gls{cpu} executes a write instruction, two primary questions arise: what happens to the cache, and what happens to lower-level memory?

\paragraph{Core Write Strategies}
\begin{itemize}
    \item \textbf{Write-Through:} Data is written to both the cache block and to the lower-level memory (main memory) simultaneously. 
    \begin{itemize}
        \item \emph{Pros:} Keeps memory consistent; easy to implement; data is safe if a crash occurs.
        \item \emph{Cons:} Generates heavy bus traffic and stalls the \gls{cpu} unless paired with a write buffer.
    \end{itemize}
    \item \textbf{Write-Back (Copy-Back):} Data is written \emph{only} to the cache block. The modified cache block is marked as ``dirty''. Main memory is only updated when that dirty block is eventually evicted from the cache to make room for a new block.
    \begin{itemize}
        \item \emph{Pros:} Significantly reduces memory bandwidth and bus traffic; fast write operations.
        \item \emph{Cons:} Complex hardware design; main memory holds stale data until eviction occurs, posing risks in multi-core or I/O direct memory access scenarios.
    \end{itemize}
\end{itemize}

\paragraph{Managing Write Misses}
When a write instruction results in a cache miss, the system must choose a policy:
\begin{itemize}
    \item \textbf{Write-Allocate:} The missing block is fetched from lower-level memory into the cache, and then the write is performed on the cache line. (Typically paired with Write-Back).
    \item \textbf{No-Write-Allocate (Write Around):} The block is modified directly in lower-level memory, bypassing the cache entirely without loading it into the cache. (Typically paired with Write-Through).
\end{itemize}

\subsection{Cache Performance and AMAT}
Cache performance is fundamentally measured in terms of \textbf{Average Memory Access Time (AMAT)}. Since memory accesses result in either a hit or a miss, the AMAT calculates the statistical average time it takes for the \gls{cpu} to retrieve data.

\begin{concept}{Average Memory Access Time (AMAT)}{amat}
    The formula for AMAT can be mathematically simplified as follows:
    \begin{align*}
        \text{AMAT} &= (\text{Hit Rate} \times \text{Hit Time}) + (\text{Miss Rate} \times \text{Miss Time}) \\
        &= (\text{Hit Rate} \times \text{Hit Time}) + \text{Miss Rate} \times (\text{Hit Time} + \text{Miss Penalty}) \\
        &= \text{Hit Time} + (\text{Miss Rate} \times \text{Miss Penalty})
    \end{align*}
\end{concept}

\noindent In modern architectures, the \textbf{Miss Penalty} is largely determined by physical hardware and technology constraints (e.g., the physical distance to main memory and bus speeds). Conversely, the \textbf{Miss Rate} can be heavily reduced through smart architectural decisions (e.g., cache size, associativity, and block size). Therefore, the primary goal in cache design is to maximize the hit rate so that the overall AMAT converges closely toward the base hit time.

\paragraph{Impact on Overall \glsfmtshort{cpu} Time}
Because the \gls{cpu} must stall and wait during a cache miss, cache performance directly impacts the total execution time of a program. The total \gls{cpu} time can be calculated by factoring in memory stall cycles:

\begin{align*}
    \text{CPU}_{\text{time}} &= (\text{CPU}_{\text{execution\_cycles}} + \text{MEM}_{\text{stall\_cycles}}) \times T_{\text{CLK}} \\
    &= \left( (IC \times \text{CPI}_{\text{exec}}) + (\text{MEM}_{\text{accesses}} \times \text{Miss Rate} \times \text{Miss Penalty}) \right) \times T_{\text{CLK}} 
\end{align*}

\noindent To make this formula easier to apply on a per-instruction basis, we can factor out the Instruction Count ($IC$). This gives us the \textbf{Memory Accesses per Instruction} ($\frac{\text{MEM}_{\text{accesses}}}{IC}$):

\begin{align*}
    \text{CPU}_{\text{time}} &= IC \times \left( \text{CPI}_{\text{exec}} + \frac{\text{MEM}_{\text{accesses}}}{IC} \times \text{Miss Rate} \times \text{Miss Penalty} \right) \times T_{\text{CLK}}
\end{align*}



\backmatter
\chapter{Appendix}

% Generate automated lists of definitions, concepts, and theorems
\tcblistof[\section]{deflist}{List of Definitions}
\tcblistof[\section]{cptlist}{List of Concepts}
\tcblistof[\section]{thmlist}{List of Theorems}

\clearpage
% Set a nice gap between the two columns
\setlength{\columnsep}{1cm}

% Span the section title across both columns
\begin{multicols}{2}[\section{List of Acronyms}]
    \raggedright
    % Disable the default glossary title so it doesn't print twice
    \renewcommand{\glossarysection}[2][]{}
    \printglossary[type=\acronymtype]
\end{multicols}

\end{document}